% Options for packages loaded elsewhere
\PassOptionsToPackage{unicode}{hyperref}
\PassOptionsToPackage{hyphens}{url}
%
\documentclass[
]{article}
\usepackage{amsmath,amssymb}
\usepackage{iftex}
\ifPDFTeX
  \usepackage[T1]{fontenc}
  \usepackage[utf8]{inputenc}
  \usepackage{textcomp} % provide euro and other symbols
\else % if luatex or xetex
  \usepackage{unicode-math} % this also loads fontspec
  \defaultfontfeatures{Scale=MatchLowercase}
  \defaultfontfeatures[\rmfamily]{Ligatures=TeX,Scale=1}
\fi
\usepackage{lmodern}
\ifPDFTeX\else
  % xetex/luatex font selection
\fi
% Use upquote if available, for straight quotes in verbatim environments
\IfFileExists{upquote.sty}{\usepackage{upquote}}{}
\IfFileExists{microtype.sty}{% use microtype if available
  \usepackage[]{microtype}
  \UseMicrotypeSet[protrusion]{basicmath} % disable protrusion for tt fonts
}{}
\makeatletter
\@ifundefined{KOMAClassName}{% if non-KOMA class
  \IfFileExists{parskip.sty}{%
    \usepackage{parskip}
  }{% else
    \setlength{\parindent}{0pt}
    \setlength{\parskip}{6pt plus 2pt minus 1pt}}
}{% if KOMA class
  \KOMAoptions{parskip=half}}
\makeatother
\usepackage{xcolor}
\setlength{\emergencystretch}{3em} % prevent overfull lines
\providecommand{\tightlist}{%
  \setlength{\itemsep}{0pt}\setlength{\parskip}{0pt}}
\setcounter{secnumdepth}{-\maxdimen} % remove section numbering
\ifLuaTeX
  \usepackage{selnolig}  % disable illegal ligatures
\fi
\IfFileExists{bookmark.sty}{\usepackage{bookmark}}{\usepackage{hyperref}}
\IfFileExists{xurl.sty}{\usepackage{xurl}}{} % add URL line breaks if available
\urlstyle{same}
\hypersetup{
  hidelinks,
  pdfcreator={LaTeX via pandoc}}

\author{}
\date{}

\begin{document}

\hypertarget{esercizi-guida-svolti}{%
\section{Esercizi Guida (Svolti)}\label{esercizi-guida-svolti}}

\textbf{Esercizio 1: Traslazione Pura}

\textbf{Problema:} L'immagine è stata spostata di +10 pixel a destra e +5 in basso. Trova l'origine del pixel di output (15, 25).

\begin{itemize}
\tightlist
\item
  Matrice \(T^{-1}\):
  \(T^{-1} = \begin{pmatrix} 1 & 0 & 0 \\ 0 & 1 & 0 \\ -10 & -5 & 1 \end{pmatrix}\)
\item
  \textbf{Calcolo:}
  ◦ \(v = 15 - 10 = \mathbf{5}\)
  ◦ \(w = 25 - 5 = \mathbf{20}\)
\item
  \textbf{Risultato:} Il colore viene preso da (5, 20)
\end{itemize}

\begin{center}\rule{0.5\linewidth}{0.5pt}\end{center}

\textbf{Esercizio 2: Scalatura (Ingradimento)}

\textbf{Problema}: L'immagine è stata raddoppiata (s = 2). Trova l'origine del pixel di input (12, 30).

\begin{itemize}
\tightlist
\item
  Matrice \(T^{-1}\):
  \(T^{-1} = \begin{pmatrix} 0.5 & 0 & 0 \\ 0 & 0.5 & 0 \\ 0 & 0 & 1 \end{pmatrix}\)
\item
  \textbf{Calcolo:}
  ◦ \(v = 12 \cdot 0.5 = \mathbf{6}\)
  ◦ \(w = 30 \cdot 0.5 = \mathbf{15}\)
\item
  \textbf{Risultato:} Il colore viene preso da (6, 15)
\end{itemize}

\begin{center}\rule{0.5\linewidth}{0.5pt}\end{center}

Esercizio 3: Trasformazione Composta

\textbf{Problema}: L'immagine è stata spostata di +6 (x) e poi ingrandita di 2. Trova l'origine del pixel (20, 10).

\begin{itemize}
\tightlist
\item
  Matrice \(T^{-1}\):
  \(T^{-1} = \begin{pmatrix} 0.5 & 0 & 0 \\ 0 & 0.5 & 0 \\ -6 & 0 & 1 \end{pmatrix}\)
\item
  \textbf{Calcolo:}
  ◦ \(v = (20 \cdot 0.5) - 6 = 10 - 6 = \mathbf{4}\)
  ◦ \(w = (10 \cdot 0.5) = \mathbf{5}\)
\item
  \textbf{Risultato:} Il colore viene preso da (4, 5)
\end{itemize}

\end{document}
